\subsection{Elliptic conventions and the exact Hodge splitting}

Throughout this subsection, $k\subset\Qb\subset\C$ is equipped with its indicated embedding. We use the degree-one motive of an elliptic curve, in the category of motives for absolute Hodge cycles generated by abelian varieties. The calculations of realizations below are also calculations of the ordinary algebraic de Rham and singular cohomology groups. Write
\[
 E/k:\quad y^2=4x^3-g_2x-g_3,\qquad
 \Delta=g_2^3-27g_3^2\ne0,
 \quad \omega=dx/y,\quad \xi=x\,dx/y.
\]
The classes $(\omega,\xi)$ form a $k$-basis of $H^1_{\dR}(E/k)$ with $F^1=k\omega$. The transcribed page 18:13 specifies a class in the additive quotient $\C/\Qb$, without choosing a unique second differential. Our scalar $h$ below is a representative relative to this explicit basis. This choice is essential when an exact numerical value, rather than algebraicity, is asserted.

Uniformize $E(\C)$ by $\C/\Lambda$, where
\[
 \Lambda=\Z\omega_1+\Z\omega_2,\quad
 \tau=\omega_2/\omega_1,\quad\operatorname{Im}\tau>0,
 \quad (x,y)=(\wp_\Lambda(z),\wp'_\Lambda(z)).
\]
The $\omega_i$ are full periods. Set $\eta_i=\zeta_\Lambda(z+\omega_i)-\zeta_\Lambda(z)$. Since $\zeta'=-\wp$, the periods of $\xi$ are $-\eta_i$. For the corresponding oriented homology basis, integration is the $\C$-linear map
\[
 I:H^1_{\dR}(E/k)\otimes_k\C
 \longrightarrow \operatorname{Hom}_{\Z}(H_1(E(\C),\Z),\C),
 \qquad I(v)(\gamma)=\int_\gamma v.
\]
With forms as rows and cycles as columns its matrix and determinant are
\begin{equation}\label{ell:periodmatrix}
 P=\begin{pmatrix}\omega_1&\omega_2\\-\eta_1&-\eta_2\end{pmatrix},
 \qquad \det P=\eta_1\omega_2-\eta_2\omega_1=2\pi i.
\end{equation}
This is Legendre's relation with full periods and the stated signs.

\begin{result}{A.4.1}{The scalar of the de Rham Hodge splitting}{Proved}
Writing $H^{0,1}=\C(\xi+h\omega)$, one has
\begin{align}
 h&=\frac{\eta_1\overline{\omega_2}-\eta_2\overline{\omega_1}}
 {\omega_1\overline{\omega_2}-\omega_2\overline{\omega_1}}\label{ell:hperiod}\\
 &=\frac{\eta_1}{\omega_1}-\frac{\pi}{\omega_1^2\operatorname{Im}\tau}
 =\frac{\pi^2}{3\omega_1^2}E_2^*(\tau).\label{ell:he2}
\end{align}
For an algebraic adapted change of basis $\omega'=r\omega$, $\xi'=s\xi+t\omega$, with $r,s\in\Qb^\times$, $t\in\Qb$, the parameter is $h'=(sh-t)/r$. In particular, the property $h\in\Qb$ is independent of the adapted algebraic basis.
\end{result}
\begin{proof}
The conjugate of the holomorphic class has periods $\overline{\omega_i}$. Write it as $a\omega+b\xi$; thus $a\omega_i-b\eta_i=\overline{\omega_i}$. Put $A=\operatorname{Im}(\overline{\omega_1}\omega_2)=|\omega_1|^2\operatorname{Im}\tau>0$. Solving the two equations by \eqref{ell:periodmatrix} gives $b=-A/\pi\ne0$ and $a/b$ equal to \eqref{ell:hperiod}. Substituting $\omega_2=\tau\omega_1$ and $\eta_2=\tau\eta_1-2\pi i/\omega_1$ gives the first expression in \eqref{ell:he2}. The identity $\eta_1\omega_1=\pi^2 E_2(\tau)/3$ gives the second. Finally,
$\xi'+h'\omega'=s(\xi+h\omega)$ if and only if $t+rh'=sh$, which proves the asserted basis transformation.
\end{proof}

Here the Fourier and Weierstrass normalizations are
\begin{align*}
 q&=e^{2\pi i\tau},& E_2&=1-24\sum_{n\ge1}\sigma_1(n)q^n,
 &E_2^*&=E_2-\frac3{\pi\operatorname{Im}\tau},\\
 E_4&=1+240\sum_{n\ge1}\sigma_3(n)q^n,
 &E_6&=1-504\sum_{n\ge1}\sigma_5(n)q^n,\\
 g_2&=\frac{4\pi^4}{3\omega_1^4}E_4,
 &g_3&=\frac{8\pi^6}{27\omega_1^6}E_6.
\end{align*}
Katz's $S(E,\omega)$ equals $-12h$. His Fourier choice $\omega_1=2\pi i$ consequently gives $S=E_2^*$. This is the normalization of \cite[Lemma 1.3.6 and (1.3.7), p.467; (1.3.10), p.468]{EllKatz1976}, with the algebraic de Rham basis in his \S1.2. Replacing full periods by half-periods requires changing Legendre's relation as well.

\subsection{The filtration field and the bigrading field}

\begin{result}{A.3.1}{The two elliptic flag fields}{Proved, using Schneider's theorem}
In the rational Betti basis dual to the cycles used above,
\[
 F^1=\C(1,\tau),\qquad H^{0,1}=\C(1,\overline\tau),
 \qquad \Q(\varphi)=\Q(\tau),\quad
 \Q(\psi)=\Q(\tau,\overline\tau).
\]
For $E/\Qb$, the filtration field has transcendence degree zero in the CM case and one in the non-CM case. The bigrading field has transcendence degree zero in the CM case; in the non-CM case its transcendence degree lies between one and two.
\end{result}
\begin{proof}
Integration of $\omega$ gives $(\omega_1,\omega_2)$, and integration of its complex conjugate gives $(\overline{\omega_1},\overline{\omega_2})$. Both first coordinates are nonzero. The affine coordinate of each line in $\mathbf P^1$ is therefore exactly its period ratio. The ordered-pair chart for the two transverse lines has coordinates $\tau,\overline\tau$, proving the field equalities. Schneider's original theorem, \cite[\S1, result IIc, p.5]{EllSchneider1937}, says that an algebraic $j(\tau)$ permits algebraic $\tau$ only when $\tau$ is imaginary quadratic. Conversely an imaginary-quadratic period ratio gives complex multiplication: multiplication by a suitable nonrational integral element of $\Q(\tau)$ preserves a finite-index sublattice and, after multiplying the element by an integer, preserves $\Z+\Z\tau$. Thus the two cases have the stated transcendence degrees. In the CM case $\overline\tau$ is also algebraic; the remaining bounds follow from the displayed field equalities.
\end{proof}

\begin{result}{A.3.2}{Curves isogenous to their complex conjugates}{Proved}
If $E/\Qb$ is isogenous to $\overline E$, then $\overline\tau\in\Q(\tau)$, so $\Q(\psi)=\Q(\varphi)$. In the non-CM case their common transcendence degree over $\Q$ is one. Moreover
\[
 h\in\Qb(\omega_1,\omega_2,\eta_1,\eta_2).
\]
This applies to every curve with a real model and also to curves isogenous to their conjugates without having a real model.
\end{result}
\begin{proof}
A holomorphic isogeny is induced on universal covers by a nonzero scalar $\lambda$ with $\lambda\Lambda\subset\overline\Lambda$. Therefore there are integers $a,b,c,d$ with $ad-bc\ne0$ such that
\[
 \lambda\omega_1=a\overline{\omega_1}+b\overline{\omega_2},
 \qquad \lambda\omega_2=c\overline{\omega_1}+d\overline{\omega_2}.
\]
It follows that
\[
 \tau=\frac{c+d\overline\tau}{a+b\overline\tau},
 \qquad \overline\tau=\frac{c-a\tau}{b\tau-d}.
\]
The last denominator cannot vanish: its vanishing would also force $c-a\tau=0$, contradicting the nonzero determinant. Conversely, a rational fractional-linear relation of this form with nonzero determinant gives commensurable lattices after scaling, hence an analytic isogeny. The isogeny descends to $\Qb$: its finite kernel consists of torsion points of the algebraic curve, its quotient is defined over $\Qb$, and a curve with the target $j$-invariant is isomorphic to the target over $\Qb$.

For an isogeny over $\Qb$, pullback of an invariant differential gives $f^*\overline\omega=\lambda\omega$ with $\lambda\in\Qb^\times$. Inverting the displayed integral matrix over $\Q$ puts $\overline{\omega_1},\overline{\omega_2}$ in $\Qb(\omega_1,\omega_2)$. Formula \eqref{ell:hperiod} proves the last inclusion. A real model gives an isomorphism to the conjugate, hence is a special case.
\end{proof}

\begin{result}{A.3.3}{An explicit nonreal example with bigrading transcendence degree one}{Proved}
The curve
\[
 E_7/\Q(i):\quad Y^2=X^3+2X^2+(\tfrac12+i)X
\]
is non-CM, has no real model, and is isogenous to its conjugate. Its bigrading field has transcendence degree one.
\end{result}
\begin{proof}
For $a=2$, $b=1/2+i$, the degree-two isogeny
\[
 (X,Y)\longmapsto\bigl(X+a+b/X,\,Y(1-b/X^2)\bigr)
\]
has target $Y'^2=X'^3-4X'^2+(2-4i)X'$. The formula follows by substitution into the target equation; the map extends at its poles to the smooth projective curves, preserves the identity, and has degree two. If $u^2=-2$, scaling $\overline E_7$ by $(X,Y)\mapsto(u^2X,u^3Y)$ gives exactly this target because $u^2\overline a=-4$ and $u^4\overline b=2-4i$.

Direct evaluation gives
\[
 j(E_7)=\frac{-56512+38144i}{25},\qquad
 N_{\Q(i)/\Q}(j(E_7))=\frac{929714176}{125}.
\]
The nonintegral norm proves that $j(E_7)$ is not an algebraic integer, whereas a CM $j$-invariant is integral; an opened proof of the latter is \cite[Theorem 12.14 and Corollaries 12.15--12.16, pp.119--120]{EllMilneMF}. Its nonreal $j$ precludes a real model because the invariant of a real equation is real and is invariant under complex isomorphism. Apply A.3.2. The short-Weierstrass invariants are $g_2=10/3-4i$, $g_3=-28/27+8i/3$. The symbolic program \texttt{check\_conjugate\_examples.py} verifies all these identities.
\end{proof}

\begin{result}{A.3.4}{An explicit curve not isogenous to its conjugate}{Proved}
Let $j_0=1/(2+i)$, $c=j_0/(j_0-1728)$, and
\[
 E_8/\Q(i):\quad y^2=4x^3-3cx-c.
\]
Then $j(E_8)=j_0$. This curve is non-CM, has no real model, and is not isogenous to $\overline E_8$.
\end{result}
\begin{proof}
Substitution in $j=1728g_2^3/(g_2^3-27g_3^2)$ gives $j_0$; the discriminant is nonzero because $c\ne0,1$. At the Gaussian prime $v=(2+i)$ one has $v(j_0)=-1$, but $v(\overline{j_0})=0$: in the residue field $\mathbf F_5$, $i=3$ and $2-i=4\ne0$. Nonintegrality gives non-CM, and nonreality gives the absence of a real model as in A.3.3.

Any isogeny between complex elliptic curves can be replaced by a cyclic isogeny between the same isomorphism classes. Indeed a kernel of type $\Z/m\oplus\Z/n$, $m\mid n$, contains $E[m]$: its $m$-torsion has order $m^2$ and equals $E[m]$. Factoring through $[m]$ leaves a cyclic kernel of order $n/m$. A cyclic $N$-isogeny between $E_8$ and its conjugate would give $\Phi_N(j_0,\overline{j_0})=0$. The modular polynomial has integer coefficients and is monic in either variable, by \cite[Theorem 6.1 and its proof, pp.82--83]{EllMilneMF}. Consequently $j_0$ would be integral over the valuation ring containing $\overline{j_0}$, contradicting $v(j_0)=-1$. For $N=1$ the same conclusion follows from equality of the $j$-invariants. The defining identities are checked symbolically in the cited program.
\end{proof}

\paragraph{Unresolved gate A.3.}
For $E$ and $\overline E$ nonisogenous, the remaining assertion is algebraic independence of $\tau,\overline\tau$. A.3.2 excludes rational fractional-linear relations in this case; its proof does not exclude polynomial relations of higher degree. For the explicit $E_8$, the bounds proved here remain
$1\le\trdeg_{\Q}\Q(\tau,\overline\tau)\le2$. No opened theorem in this run settles this gate.

\subsection{CM algebraicity, its exact value, and the CM period field}

\begin{result}{A.4.2}{Algebraic Hodge splitting for CM curves}{Known, with derivation}
For a CM elliptic curve $E/k$ in the stated Weierstrass normalization, $h\in k$. In particular its Hodge splitting is defined over $k$.
\end{result}
\begin{proof}
First choose a nonreal CM endomorphism $u$, defined over a finite extension of $k$. On the holomorphic and antiholomorphic lines its eigenvalues are $u$ and $\overline u$, respectively. In the de Rham basis its action is
\[
 u^*\omega=u\omega,\qquad u^*\xi=a\omega+\overline u\,\xi,
 \qquad a\in\Qb.
\]
Applying it to $\xi+h\omega$ gives $a+uh=\overline u h$, hence $h=a/(\overline u-u)\in\Qb$. This is the eigenspace mechanism of \cite[Lemma 4.0.7, p.501]{EllKatz1976}. To descend to the coefficient field itself, Masser's Theorem A1 gives $\Psi(\tau)\in\Q(j(\tau))$ for his function
$\Psi=3E_4E_2^*/(2E_6)$ at an imaginary-quadratic point outside the orbit of $i$; see \cite[Appendix One, Theorem A1 and (100), p.114]{EllMasser1975}. If $j\ne0,1728$, the explicit normalizations give $h=(g_3/g_2)\Psi(\tau)\in k$. At $j=0,1728$, the modular stabilizers force $E_2^*=0$ and therefore $h=0\in k$. This also handles the two exceptional denominators in the displayed formula.
\end{proof}

\begin{result}{A.4.3}{The full CM period field}{Known, with exact field calculation}
For a CM elliptic curve over $\Qb$, let
\[
 L=\Qb(\omega_1,\omega_2,\eta_1,\eta_2).
\]
Then
\[
 L=\Qb(\omega_1,\pi/\omega_1)=\Qb(\omega_1,\pi),
 \qquad \trdeg_{\Qb}L=2.
\]
For the degree-one motive, the CM motivic torus has dimension two; this proves the period-conjecture dimension in this case. The same transcendence degree holds after adjoining the periods of the conjugate curve.
\end{result}
\begin{proof}
The numbers $\tau$ and $\operatorname{Im}\tau$ are algebraic and nonzero in the CM case, and A.4.2 gives algebraic $h$. Formula \eqref{ell:he2} gives
\[
 \omega_2=\tau\omega_1,\qquad
 \eta_1=h\omega_1+\frac{\pi}{\omega_1\operatorname{Im}\tau},
 \qquad \eta_2=\tau\eta_1-\frac{2\pi i}{\omega_1}.
\]
Thus $L\subset\Qb(\omega_1,\pi/\omega_1)$. Legendre's relation places $\pi$ in $L$, proving the reverse inclusion. Chudnovsky's theorem gives transcendence degree at least two for the four periods and quasiperiods of an elliptic curve with algebraic invariants; we use the original author's \cite[\S2A, Theorem 1, p.343, with the algebraic-invariant convention on p.339]{EllChudnovsky1978}. The displayed field has at most two generators, so equality follows. This also proves algebraic independence of $\pi$ and $\omega_1$.

The CM action confines the Mumford--Tate group to the two-dimensional torus $\operatorname{Res}_{K/\Q}\Gm$, where $K$ is the CM field. Over $\C$ the Hodge cocharacter and its Galois conjugate are the two coordinate cocharacters, so they generate that torus. In the category of abelian motives for absolute Hodge cycles this is the motivic group, by the absolute-Hodge theorem used for the category in the paper.

For completeness, the conjugate CM lattice is commensurable after scaling with the original: both normalized lattices span the same imaginary-quadratic field over $\Q$. The resulting isogeny descends to $\Qb$ by the kernel argument in A.3.2. Its matrix on de Rham cohomology is invertible over $\Qb$, and its matrix on rational homology is invertible over $\Q$. Integration therefore expresses all periods and quasiperiods of the conjugate curve as $\Qb$-linear combinations of those of $E$. The compositum has the same field $L$. For any algebraic base field $k$ over which the data are defined, algebraic extension to $\Qb$ does not change these transcendence degrees.
\end{proof}

\begin{result}{A.4.4}{An exact nonzero CM splitting}{Proved}
For $E:y^2=4x^3-30x-28$ one has $j=8000$, CM by $\Q(\sqrt{-2})$, and $h=1/2$ in the basis $(dx/y,x\,dx/y)$.
\end{result}
\begin{proof}
Put $X=x+2$, $Y=y/2$. The equation becomes $E_0:Y^2=X^3-6X^2+9X/2$. The formulas
\[
 X'=X-6+\frac9{2X},\qquad Y'=Y\left(1-\frac9{2X^2}\right)
\]
give a normalized degree-two isogeny $\varphi:E_0\to E_1$, where $E_1:Y'^2=X'^3+12X'^2+18X'$. Substitution verifies the target equation, and pullback sends $dX'/(2Y')$ to $dX/(2Y)$. The induced function-field extension has degree two, as is also seen from the quadratic equation for $X$ in terms of $X'$. The rational map extends to the smooth projective curves and carries the origin to the origin.

For $u=\sqrt{-2}$, the map $s:E_0\to E_1$, $(X,Y)\mapsto(u^2X,u^3Y)$, is an isomorphism. Hence $\alpha=s^{-1}\varphi$ is a degree-two endomorphism with $\alpha^*\omega=u\omega$. Its nonrational eigenvalue proves CM. Moreover
\[
 d(Y/X)=\left(X-\frac9{2X}\right)\omega,
 \qquad \alpha^*\xi/u=(-X'/2-2)\omega
 =(-X+1)\omega+\tfrac12d(Y/X).
\]
Because $\xi=(X-2)\omega$, this gives $\alpha^*\xi=-u\xi-u\omega$ in de Rham cohomology. The eigenspace calculation in A.4.2 yields $h=(-u)/(-2u)=1/2$. The discriminant is $5832$, giving $j=1728\cdot30^3/5832=8000$. The program \texttt{verify\_cm\_isogeny.py} checks every rational-function identity exactly.
\end{proof}

\begin{result}{A.4.5}{Masser's explicit CM formula in these conventions}{Known, with normalization checked}
For imaginary-quadratic $\tau$ with $j=j(\tau)\ne0,1728$, define $\Psi=3E_4E_2^*/(2E_6)$. Then
\[
 h=(g_3/g_2)\Psi,\qquad
 \Psi=9j\gamma+\frac{3(7j-6912)}{2(j-1728)},
\]
where $\gamma\in\Q(j)$ is the modular-polynomial expression specified below. At $j=0,1728$, one has $h=0$.
\end{result}
\begin{proof}[Derivation and exact locator]
Choose a primitive equation $C\tau^2+B\tau+A=0$, $C>0$, and set
\[
 \mu=\begin{cases}
 \begin{pmatrix}-B&-2A\\2C&B\end{pmatrix},&B\text{ odd},\\[2mm]
 \begin{pmatrix}-B/2&-A\\C&B/2\end{pmatrix},&B\text{ even},
 \end{cases}\qquad D=\det\mu.
\]
Expand $\Phi_D(j+X,j+Y)=\sum\beta_{mn}X^mY^n$. In the ordinary case,
\[
 \beta=\beta_{10}=\beta_{01}\ne0,\qquad
 \gamma=(\beta_{20}-\beta_{11}+\beta_{02})/\beta.
\]
In Masser's exceptional case, $B$ is odd and $D=3d^2$ for an integer $d>1$, and
\[
 \beta=\beta_{30}=\beta_{03}\ne0,\quad \beta_{21}=\beta_{12}=0,
 \quad \gamma=(\beta_{40}-\beta_{31}+\beta_{22}-\beta_{13}+\beta_{04})/\beta.
\]
These are \cite[Lemmas A1--A2, pp.116--118, and (105)--(106), pp.118--119]{EllMasser1975}. The equality for $h$ follows by dividing the exact expressions for $g_3,g_2$ above. The application of (106) is restricted here to $j\ne0,1728$: its displayed derivation divides by $j'(\tau)$. At the elliptic points one instead uses $E_2^*=0$. Naive substitution of $j=0$ in the rational expression would give a spurious value. The table on Masser's p.121 gives $\Psi(i\sqrt2)=15/28$, and $(28/30)(15/28)=1/2$, independently checking the normalization in A.4.4.
\end{proof}

\paragraph{Two transcription hazards in the later literature.}
For numerical formulas, use Masser's original (100). The abbreviated description of his function in \cite[p.774]{EllSpence2018} omits its factor $3/2$; the field-inclusion argument is unaffected. On that same printed page, the exceptional expression displays the denominator $\beta_{01}$, whereas Masser's exceptional denominator is $\beta_{03}$ and $\beta_{01}$ vanishes. Both printed pages were visually checked; this is not an OCR substitution. The statements used here are taken from Masser's original formula.

\subsection{The zero value and isogeny transport}

\begin{result}{A.4.6}{Complete classification of the zero scalar}{Proved}
For any complex elliptic curve in the short-Weierstrass normalization fixed above,
\[
 h=0\quad\Longleftrightarrow\quad j(E)\in\{0,1728\}.
\]
\end{result}
\begin{proof}
By \eqref{ell:he2}, it suffices to locate the zeros of $E_2^*$. Its weight-two transformation law reduces the question to $|x|\le1/2$, $|\tau|\ge1$, for $\tau=x+iy$. Thus $y\ge\sqrt3/2$ and $r=e^{-2\pi y}<1/200$. For $0<x<1/2$, use $\sigma_1(n)\le n^2$ and $|\sin(n\theta)|\le n\sin\theta$ for $0<\theta<\pi$ to obtain
\[
 \left|\sum_{n\ge2}\sigma_1(n)r^n\sin(2\pi nx)\right|
 \le\sin(2\pi x)\sum_{n\ge2}n^3r^n
 <r\sin(2\pi x).
\]
Indeed
\[
 S_3(r)=\sum_{n\ge1}n^3r^n=\frac{r(1+4r+r^2)}{(1-r)^4},
 \qquad S_3(1/200)/(1/200)<2.
\]
The first Fourier term therefore gives a nonzero imaginary part. Reflection deals with $-1/2<x<0$. A zero must lie on $x=0$ or $x=\pm1/2$.

For $x=0$, $y\ge1$ and
\[
 \frac d{dy}E_2^*(iy)=48\pi\sum_{n\ge1}n\sigma_1(n)e^{-2\pi ny}
 +\frac3{\pi y^2}>0.
\]
For $x=1/2$, $y\ge y_0=\sqrt3/2$, and $y^2e^{-2\pi ny}$ decreases on this ray. Consequently
\[
 y^2\frac d{dy}E_2^*(1/2+iy)
 \ge\frac3\pi-36\pi S_3(e^{-\pi\sqrt3})
 >\frac3\pi-\frac\pi5>0.
\]
Here $S_3(1/200)=8160200/1568239201<1/180$ and $\pi^2<15$. Modular stabilizers force zero at $i$ and at $\rho+1$, where $\rho=e^{2\pi i/3}$. Strict monotonicity proves uniqueness on each ray; reflection and translation finish the fundamental domain. The orbits have $j=1728$ and $j=0$.

The analytic lemma and proof are also in the opened primary source \cite[Lemma 8.1, printed pp.25--26]{EllSebbarBesrour2026}. Every estimate just used was independently checked with rational bounds by \texttt{check\_zero\_bounds.py}. This assertion concerns the scalar in the fixed short-Weierstrass basis. Its class in $\C/\Qb$ vanishes for all CM curves, including A.4.4 where the scalar is nonzero.
\end{proof}

\begin{result}{A.4.7}{Algebraic isogeny transport of the splitting scalar}{Proved}
Let $f:E\to E'$ be a degree-$n$ isogeny over $\Qb$ between short-Weierstrass models, and write
\[
 f^*\omega'=\lambda\omega,\qquad f^*\xi'=a\omega+b\xi.
\]
Then $\lambda,a,b\in\Qb$, $\lambda b=n$, and
\[
 h_{E'}=\frac{n}{\lambda^2}h_E-\frac a\lambda.
\]
Hence algebraicity of the splitting is constant on an algebraic isogeny class.
\end{result}
\begin{proof}
Pullback on algebraic de Rham cohomology has algebraic coefficients and preserves $F^1$. The cup-product pairing is normalized by $(\omega,\xi)_{\dR}=1$, as in Katz \S1.2. Pullback multiplies this pairing by $n$, so its determinant is $\lambda b=n$ and both factors are nonzero. Preservation of $H^{0,1}$ gives
\[
 a\omega+b\xi+\lambda h_{E'}\omega=b(\xi+h_E\omega).
\]
Solving for $h_{E'}$ proves the formula and the assertion.
\end{proof}

\subsection{The almost-holomorphic formulation and the exact remaining gate}

\begin{result}{A.4.8}{Translation into the function $\chi^*$}{Proved}
For $j\ne0,1728$, define
\[
 \chi^*(\tau)=\frac{1728E_2^*(\tau)E_4(\tau)E_6(\tau)}
 {E_4(\tau)^3-E_6(\tau)^2}.
\]
Then
\[
 \chi^*=\tfrac23(j-1728)\Psi,\qquad
 h=\frac{3g_3}{2g_2(j-1728)}\chi^*.
\]
For algebraic $g_2,g_3\ne0$, algebraicity of $h$ is equivalent to algebraicity of $\chi^*(\tau)$.
\end{result}
\begin{proof}
Substitute $j-1728=1728E_6^2/(E_4^3-E_6^2)$ and the definition of $\Psi$. The coefficient in the second equality is a nonzero algebraic number under the stated hypotheses, proving the equivalence.
\end{proof}

Spence's \cite[Proposition 5.1, p.773]{EllSpence2018} proves $\Q(\chi^*(\tau))\subseteq\Q(j(\tau))$ at quadratic $\tau$. His Theorem 5.8, p.776, proves finiteness of maximal $(j,\chi^*)$-special subvarieties. Definition 3.6, pp.755--756, defines their constant factors using quadratic $\tau$. Applied to an arbitrary algebraic singleton, the finiteness theorem allows either a special singleton or no special subvariety at all. This specialization therefore does not prove that every algebraic pair $(j(\tau),\chi^*(\tau))$ is CM. The calculation states exactly what that application yields; it does not exclude a further argument extending the methods.

\paragraph{Unresolved gate A.4(c).}
The converse in Katz's \cite[Question 4.0.8, p.501]{EllKatz1976} is not established by the results above. Its concrete real test case here is
\[
 E_* /\Q:\quad y^2=4x^3-4x-1,
 \qquad \Delta=37,\quad j=110592/37.
\]
Nonintegrality of $j$ proves non-CM. The cubic has three distinct real roots, so choose $\omega_1>0$ real and $\omega_2\in i\R_{>0}$. Reality of the Weierstrass series gives $\eta_1\in\R$, $\eta_2\in i\R$. Formula \eqref{ell:hperiod} becomes
\[
 h_* =\frac{\eta_1\omega_2+\eta_2\omega_1}{2\omega_1\omega_2}.
\]
For any $a\in\Qb$, there are exact equivalences
\begin{equation}\label{ell:exactgate}
 \begin{split}
 h_*=a
 &\quad\Longleftrightarrow\quad
 \eta_1\omega_2+\eta_2\omega_1-2a\omega_1\omega_2=0\\
 &\quad\Longleftrightarrow\quad (\eta_1-a\omega_1)\omega_2=\pi i.
 \end{split}
\end{equation}
The second equivalence follows from Legendre's relation with the signs in \eqref{ell:periodmatrix}. The unresolved assertion is that \eqref{ell:exactgate} fails for every algebraic $a$. A.4.6 excludes $a=0$; it does not exclude the other algebraic values. The computed approximation is
\[
 h_*=0.113219334661746861539065791000647532558647974361219962657816837416979\ldots.
\]
No algebraicity inference is drawn from it. For curves without a real model, formula \eqref{ell:hperiod} is the exact obstruction with the conjugate periods retained; A.3.2 proves their elimination when the curve is isogenous to its conjugate.

\subsection{Tests of the classical transcendence theorems}

Chudnovsky's \cite[\S2A, Theorems 1 and 3, p.343]{EllChudnovsky1978} gives, respectively, transcendence degree at least two for the four periods and quasiperiods, and algebraic independence of $\pi/\omega$ and $\eta/\omega$. Here is an explicit algebraic test of the proposed contradiction from the first assertion and linear period independence. Write
\[
 x=\omega_1,\quad y=\omega_2,\quad
 u=\eta_1-a\omega_1,\quad v=\eta_2-a\omega_2.
\]
The putative relation is $uy+vx=0$. On the open set $xyu\ne0$, the algebraic variety $v=-uy/x$ has function field $\Qb(x,y,u)$ of transcendence degree three, and its four coordinate functions are $\Qb$-linearly independent. Indeed multiplying a proposed constant linear relation by $x$ gives a linear combination of the distinct monomials $x^2,xy,xu,uy$. Its determinant $uy-vx=2uy$ is nonzero. This is an algebraic compatibility witness for the displayed relation and those lower bounds, not a tuple of genuine periods and not a counterexample to a stronger period theorem.

Nesterenko's \cite[Theorem 1, p.1320]{EllNesterenko1996} gives transcendence degree at least three for $q,E_2(q),E_4(q),E_6(q)$ when $0<|q|<1$. His Corollary 3, p.1321, gives algebraic independence of
\[
 q,\qquad \omega_1/\pi,\qquad \eta_1/\pi
\]
for algebraic invariants. With $t=\omega_1/\pi$, the equality $h=a\in\Qb$ in \eqref{ell:he2} gives
\[
 \eta_1/\pi=at+\frac1{t\pi\operatorname{Im}\tau},
 \quad
 \Qb(q,t,\eta_1/\pi)=\Qb(q,t,\pi\operatorname{Im}\tau).
\]
Nesterenko's theorem then gives transcendence degree three for the generators on the right. It does not contradict this field equality. This checks the specialization of the exact theorem, with its normalizations retained.

The original Schneider source was obtained as the G\"ottingen scan and read at \S1, result IIc, p.5. The complete original Chudnovsky congress contribution was obtained from the official IMU volume and read at pp.339 and 343. The 2026 theorem of \cite[Theorem 1]{EllSenthilKumar2026} requires $u_1,u_2,\omega$ to be $\Q$-linearly independent and $(\Z u_1+\Z u_2)\cap\Lambda=0$; its at-least-two lower bound cannot be specialized by taking the two extra logarithms in the period lattice. Waldschmidt's \cite[\S3.3, p.31]{EllWaldschmidt2004} explicitly records Katz's question and points to Bertrand's nonholomorphic conjecture. This supplies historical status evidence, not proof of its present status. The referenced Bertrand chapter's metadata was opened, but its full text was not obtained and no unseen theorem from it is invoked.

\subsection{Reproducibility of the elliptic computations}

The exact-input programs accompany this section. The isogeny and conjugate examples use symbolic rational-function arithmetic. The zero estimates use rational inequalities, including $S_3(1/200)<1/180$ and a rational Taylor lower bound proving $e^{\pi\sqrt3}>200$. The numerical program evaluates $(g_2,g_3)=(4,0),(30,28),(4,1)$ at 80 and 160 decimal digits.

For real roots $e_1>e_2>e_3$, put $d=e_1-e_3$, $m=(e_2-e_3)/d$. It independently computes
\begin{align*}
 \omega_1&=2K(m)/\sqrt d,&
 \omega_2&=2iK(1-m)/\sqrt d,\\
 \eta_1&=2\sqrt d\,\mathrm E(m)-e_1\omega_1,&
 \eta_2&=-2i\sqrt d\,\mathrm E(1-m)-e_3\omega_2.
\end{align*}
Here $K$ and $\mathrm E$ are the complete elliptic integrals. In particular $\eta_2$ is not defined by forcing Legendre's relation. The program compares these values with the Fourier formula for $\eta_1$, verifies Legendre's relation, compares both formulas for $h$, and reconstructs $g_2,g_3$. Residuals in the 160-digit run are of order $10^{-160}$.

For a Fourier sum with divisor weight $s$, coefficient $c$, $N$ retained terms, $r=|q|$, and $R=r((N+2)/(N+1))^{s+1}<1$, the omitted tail is bounded by
\[
 |c|\frac{(N+1)^{s+1}r^{N+1}}{1-R},
\]
using $\sigma_s(n)\le n^{s+1}$. This bound is rigorous for exact $r$. The implementation evaluates the integrals and $r$ in high-precision floating arithmetic, so the output is a consistency check rather than a full interval certificate for rounding error. The programs, inputs, precision, and saved outputs are all included; no transcendence conclusion rests on numerical recognition.
